\chapter{Diseño experimental}
\label{chap:diseno}

El diseño experimental se organiza alrededor de dos decisiones. La primera es clínica: el sistema se plantea como apoyo de triaje para sospecha de Brugada, no como diagnóstico autónomo. La segunda es metodológica: se comparan dos formas de usar pocos ejemplos etiquetados. La CNN aprende modificando sus pesos; el VLM previsto aprenderá en contexto, sin ajuste de parámetros.

\begin{figure}[H]
  \centering
  \safeincludegraphics[width=0.96\textwidth]{results/fig_triaje_clinico.pdf}
  \caption{Orientación clínica del prototipo: una alerta prioriza revisión y posible circuito diagnóstico, pero no decide el diagnóstico.}
  \label{fig:triaje_clinico}
\end{figure}

\section{Pregunta experimental}

La pregunta completa del TFG es cuándo renta entrenar un modelo visual específico frente a utilizar un modelo de visión y lenguaje con aprendizaje en contexto. Como la campaña VLM no dispone todavía de inferencias auditadas, esta memoria cierra el primer brazo de la comparación: la rama CNN.

El brazo CNN debe responder si el enfoque supervisado clásico es suficiente con tan pocos datos. Para ello se evalúa primero en un dominio controlado, después en Brugada HUCA y finalmente con adaptación de dominio. Si la CNN falla en el dominio real incluso tras estas técnicas, el resultado no es accesorio: es la evidencia que justifica ejecutar la comparación VLM/ICL.

El estudio se cierra con tres familias de experimentos:

\begin{itemize}
  \item CNN ResNet18 sobre el conjunto sintético QRS/ST con etiqueta derivada de patrón Brugada operativo.
  \item CNN ResNet18 sobre Brugada HUCA con etiqueta de caso confirmado frente a control.
  \item Adaptación de dominio desde sintético a HUCA mediante aprendizaje autosupervisado, CORAL, MMD y DANN.
\end{itemize}

Los modelos abiertos de visión y lenguaje quedan fuera de la evidencia empírica final, pero no fuera del diseño. La infraestructura, los mensajes, los controles y la validación de salida se documentan porque son el segundo brazo de la comparación. No se reportan métricas VLM porque no existen predicciones completas auditadas.

\section{Canalización experimental}

La figura \ref{fig:pipeline_experimental} resume la canalización. Las imágenes sintéticas y reales se convierten al mismo formato visual de entrada. La ruta CNN entrena modelos con presupuestos reducidos, evalúa por paciente, agrega predicciones por derivación y produce tablas, curvas, matrices de confusión y Grad CAM. La ruta VLM prevista reutiliza las mismas imágenes y presupuestos \(K\), pero introduce los ejemplos como demostraciones dentro del prompt.

\begin{figure}[H]
  \centering
  \safeincludegraphics[width=0.96\textwidth]{results/fig_pipeline_experimental.pdf}
  \caption{Diseño experimental: CNN como línea base cerrada y VLM/ICL como comparador central pendiente.}
  \label{fig:pipeline_experimental}
\end{figure}

\section{Unidades de análisis}

La unidad de separación es siempre el paciente. En el conjunto sintético, cada paciente simulado contiene tres imágenes correspondientes a V1, V2 y V3. En Brugada HUCA, cada paciente clínico aporta las ventanas renderizadas de esas derivaciones precordiales derechas. El modelo produce una probabilidad por imagen y la decisión de paciente se obtiene promediando las probabilidades de sus derivaciones antes de aplicar el umbral.

\begin{table}[H]
  \centering
  \caption{Conjuntos utilizados en los experimentos cerrados.}
  \label{tab:datasets_cerrados}
  \small
  \begin{tabularx}{\textwidth}{p{3cm}rrrrX}
    \toprule
    Conjunto & Pacientes & Imágenes & Positivos & Negativos & Tarea \\
    \midrule
    Sintético QRS/ST & 100 & 300 & 20 & 80 & Detección derivada de patrón Brugada operativo: RBBB, elevación ST e inversión T presentes \\
    Brugada HUCA & 317 & 951 & 116 & 201 & Caso Brugada confirmado frente a control, excluyendo categorías limítrofes o no equivalentes \\
    \bottomrule
  \end{tabularx}
\end{table}

La diferencia entre ambas tareas es deliberada y debe conservarse al interpretar los resultados. En el simulador se evalúa composición morfológica derivada de reglas. En HUCA se evalúa una etiqueta clínica binaria disponible en los metadatos. Por tanto, el salto sintético-real no mide únicamente cambio de renderizado, sino también un cambio de definición de etiqueta.

\section{Presupuestos de pocos ejemplos}

Los presupuestos de entrenamiento son \(K=\{2,4,8,16,32\}\). Para cada \(K\) positivo se ejecutan tres semillas: 42, 123 y 2026. La selección de ejemplos se realiza dentro de cada pliegue, sin utilizar el paciente reservado para test. El objetivo no es maximizar rendimiento con todos los datos, sino medir cómo evoluciona cada estrategia cuando las etiquetas son escasas.

\begin{figure}[H]
  \centering
  \safeincludegraphics[width=0.94\textwidth]{results/fig_k_budget_split.pdf}
  \caption{Validación leave-one-out por paciente y selección de presupuestos \(K\).}
  \label{fig:k_budget_split}
\end{figure}

En cada ejecución se conserva el identificador de pliegue, la semilla, los pacientes usados como ejemplos, las predicciones por imagen y las métricas agregadas. El informe de auditoría final verifica 1500 predicciones de paciente en el simulador y 4755 predicciones de paciente en HUCA para la CNN base.

\section{Modelo CNN evaluado}

El modelo supervisado final es ResNet18 con pesos iniciales de ImageNet y una cabeza binaria. Aunque el simulador original contiene tres hallazgos morfológicos, la campaña cerrada utiliza una etiqueta derivada de patrón Brugada operativo: positiva cuando las condiciones necesarias están simultáneamente presentes y negativa cuando el paciente simulado es normal o incompleto. Esta reducción facilita la comparación con HUCA, donde la etiqueta disponible es binaria.

El entrenamiento utiliza \texttt{BCEWithLogitsLoss}, ponderación automática de la clase positiva cuando procede, AdamW y selección interna de umbral sobre el subconjunto de validación del pliegue. La métrica principal es la exactitud equilibrada porque la prevalencia positiva es desigual y porque interesa medir sensibilidad y especificidad de forma simétrica.

\section{Adaptación de dominio}

La rama de adaptación de dominio se revisa y ejecuta como extensión del modelo base. Sus métodos son:

\begin{itemize}
  \item \textbf{SSL}: preentrenamiento autosupervisado sobre imágenes sin etiqueta antes del ajuste supervisado.
  \item \textbf{CORAL}: penalización para acercar covarianzas de características entre origen sintético y destino real.
  \item \textbf{MMD}: penalización de discrepancia entre distribuciones de características de ambos dominios.
  \item \textbf{DANN}: aprendizaje adversario con inversión de gradiente para reducir información específica del dominio \cite{ganin2016}.
\end{itemize}

Estas variantes se evalúan en \(K=16\) y \(K=32\), los presupuestos con mayor interés tras observar la curva del modelo base. La comparación se realiza contra la CNN HUCA sin adaptación en los mismos valores de \(K\). Si estas variantes no recuperan rendimiento suficiente, el hallazgo refuerza la necesidad de probar una estrategia distinta, no solo más regularización de la CNN.

\begin{figure}[H]
  \centering
  \safeincludegraphics[width=0.96\textwidth]{results/fig_domain_adaptation_methods.pdf}
  \caption{Métodos de adaptación de dominio revisados y ejecutados sobre la rama experimental.}
  \label{fig:domain_adaptation_methods}
\end{figure}

\section{Métricas}

Se reportan exactitud, exactitud equilibrada, F1, precisión, sensibilidad, especificidad, área ROC y precisión media. Las tablas principales muestran la media y la desviación estándar sobre tres semillas. Además, para cada configuración se guardan matrices de confusión agregadas, curvas ROC/PR y predicciones por paciente.

La exactitud equilibrada se define como:

\[
  \mathrm{BA}=\frac{1}{2}\left(\frac{TP}{TP+FN}+\frac{TN}{TN+FP}\right).
\]

Esta métrica es especialmente importante en HUCA, donde un modelo puede obtener sensibilidad moderada a costa de marcar demasiados controles como positivos. En una herramienta de triaje ese comportamiento puede ser tolerable solo si se declara como alerta preliminar y se acompaña de revisión clínica. En la comparación futura con VLM/ICL se usará el mismo criterio principal para evitar que el cambio de familia de modelo cambie la regla de evaluación.

\section{Estado de cierre}

\begin{table}[H]
  \centering
  \caption{Estado de los bloques experimentales al cierre de la memoria.}
  \label{tab:estado_experimentos}
  \small
  \begin{tabularx}{\textwidth}{p{3.6cm}p{2.6cm}X}
    \toprule
    Bloque & Estado & Evidencia disponible \\
    \midrule
    Simulador QRS/ST & Completo & Datos, pliegues LOOCV, entrenamiento CNN, tablas, curvas, matriz de confusión y Grad CAM \\
    Brugada HUCA & Completo & Datos procesados, pliegues por paciente, entrenamiento CNN, comparación sintético-real y Grad CAM \\
    Adaptación de dominio & Completo & SSL, CORAL, MMD y DANN ejecutados para \(K=16\) y \(K=32\), con comparación agregada \\
    VLM e ICL & TODO documentado & Segundo brazo de la comparación: código, mensajes, controles y contrato de evaluación preparados; faltan inferencias completas y métricas auditadas \\
    Auditoría final & Completa para CNN & \texttt{final\_cnn\_package\_ready: OK}; VLM marcado explícitamente como TODO \\
    \bottomrule
  \end{tabularx}
\end{table}

Este estado evita mezclar resultados cerrados con trabajo futuro. Las conclusiones cuantitativas de esta memoria se apoyan únicamente en CNN y adaptación de dominio, pero la interpretación metodológica apunta directamente a la siguiente fase: probar si VLM/ICL mejora la relación entre rendimiento, coste y cantidad de datos etiquetados.
